\section{Polynomial Evaluation with $\lfloor n/2\rfloor+1$ Multiplications and Rational Preprocessing}
\label{sec:model}

We work over a field $\F$.
An \emph{arithmetic circuit} (or straight-line program) takes inputs $(x,\theta)\in\F\times\F^d$ and outputs a value in $\F$ using additions, subtractions, and multiplications.
We measure the cost of a circuit by the number of multiplications.

Fix $n\ge 1$.
An \emph{evaluation scheme} of (target) degree $n$ is a circuit $E(x;\theta)$ such that for every $\theta\in\F^d$, the function $x\mapsto E(x;\theta)$ is a polynomial of degree at most $n$.
Let $\mathsf{coeff}(\theta)\in\F^n$ denote the coefficient vector $(c_0,\dots,c_{n-1})$ of that polynomial, where we write
\[
E(x;\theta)=c_0 + c_1 x + \dots + c_{n-1}x^{n-1} + x^n.
\]
Thus $\mathsf{coeff}:\F^d\to\F^n$ is a polynomial map.

\paragraph{Rational preprocessing/decoding.}
We say $E$ has \emph{rational preprocessing} (or \emph{rational decoding}) if there is a rational map
\[
\mathsf{pre}:\F^n \dashrightarrow \F^d
\]
such that for every coefficient vector $c$ in the domain of $\mathsf{pre}$ we have
\[
E(x;\mathsf{pre}(c)) \equiv c_0 + c_1 x + \dots + c_{n-1}x^{n-1} + x^n
\]
as polynomials in $x$.
Equivalently, $\mathsf{pre}$ is a rational right-inverse of $\mathsf{coeff}$ on a Zariski-open set, so $\mathsf{coeff}$ is generically injective (birational onto its image).

\paragraph{Hashing viewpoint.}
If we sample $\theta$ at random and evaluate $E(x;\theta)$ at an input $x$, we obtain a polynomial-hash family.
When $\mathsf{coeff}$ is (generically) injective, distinct parameters induce distinct polynomials, which is exactly what is needed to interpret preprocessing as selecting a random polynomial from a well-defined family.

\subsection{Main Result}
\label{sec:main-theorem}

Our main contribution is an evaluation scheme with essentially the minimum possible number of multiplications while still admitting rational preprocessing.

\begin{theorem}
\label{thm:main}
For every $n\ge 1$ and every field $\F$ of characteristic $\neq 2$, there is an evaluation scheme for monic degree-$n$ polynomials over $\F$ that uses $\lfloor n/2\rfloor+1$ multiplications and admits rational preprocessing.
\end{theorem}

\begin{remark}[Characteristic 2]
\label{rem:char2}
In characteristic~2, some of the decoding steps in our characteristic $\neq 2$ construction rely on dividing by even integers (e.g.\ $2$), which is undefined.
Over perfect fields of characteristic~2 (in particular finite fields $\F_{2^w}$), there is an alternative decoding strategy that replaces those divisions by (i) Lucas’ theorem for binomial coefficients modulo~$2$ and (ii) Frobenius inverses (unique $2^t$-th roots).
This makes the affected decoding steps well-defined in characteristic~$2$ (with preprocessing allowed to use Frobenius roots).
A complete, unified characteristic~$2$ construction matching \cref{thm:main} remains open.
\end{remark}

\subsection{Intuition and Construction Outline}
\label{sec:main-theorem:outline}

This section gives the high-level intuition behind the construction; the full details appear in \cref{appendix:constructions}.
The proof consists of making the invariants below precise and checking that every step preserves rational decoding.

\paragraph{From rational preprocessing to decoding coefficients.}
Fix a circuit family $E(x;\theta)$ and write $c=\mathsf{coeff}(\theta)\in\F^n$ for the output coefficient vector as in \cref{sec:model}.
Rational preprocessing asks for a rational right-inverse $\mathsf{pre}$ of $\mathsf{coeff}$ on a dense set, i.e.\ that we can (generically) recover the parameters $\theta$ from the coefficients of the resulting polynomial.
Thus, in the construction we repeatedly build polynomial families whose parameters can be \emph{decoded} from their coefficients by simple triangular relationships.

\paragraph{A warm-up: degree $3$ with two multiplications.}
Let $n=3$.
Compute $H_2(x)=x^2$ using one multiplication, and then define
\[
   Q_3(x;\alpha_0,\alpha_1,\alpha_2) \;=\; (x+\alpha_2)(H_2(x)+\alpha_1)+\alpha_0.
\]
This uses one additional multiplication, so two in total, matching $\lfloor 3/2\rfloor+1$.
Expanding $H_2(x)=x^2$ gives
\[
   Q_3(x)=x^3 + \alpha_2 x^2 + \alpha_1 x + (\alpha_1\alpha_2+\alpha_0),
\]
so the coefficient map is (almost) triangular:
$\alpha_2=\coeff{Q_3}{2}$ and $\alpha_1=\coeff{Q_3}{1}$ are read off directly, and then $\alpha_0=\coeff{Q_3}{0}-\alpha_1\alpha_2$.
This illustrates the general strategy: we use each multiplication to introduce (roughly) two new degrees of freedom, while keeping decoding algebraic and explicit.

\paragraph{Splitting viewpoint (``splittable pairs'').}
For larger degrees, it is convenient to construct the output polynomial as
\[
   P(x) \;=\; x\cdot T^{(1)}(x) + T^{(2)}(x),
\]
where $T^{(1)}$ and $T^{(2)}$ are degree-$(n-1)$ polynomials computed with shared intermediate values.
The final multiplication by $x$ contributes the ``$+1$'' in the multiplication count.
We call $(T^{(1)},T^{(2)})$ a \emph{splittable pair} (for degree $n$) if the coefficients of $P(x)=xT^{(1)}(x)+T^{(2)}(x)$ generically determine all underlying parameters (equivalently: the induced coefficient map is birational onto its image).
The core of the construction is to build splittable pairs for all target degrees while paying only about one multiplication per two coefficients.

\paragraph{Keeping decoding stable: coefficient windows and ``compatible pairs''.}
The main technical problem is that multiplication mixes coefficients.
To control this, we maintain an invariant of the following form.
We consider a pair $(P^{(1)},P^{(2)})$ and its combination $\Phi(x)\coloneqq xP^{(1)}(x)+P^{(2)}(x)$.
We ensure that there is a \emph{window} of coefficient positions (typically a top-degree range) from which one can recover the parameters introduced so far and also recover the pair $(P^{(1)},P^{(2)})$ itself.
Pairs with this property are called \emph{compatible pairs} in the long version.

Crucially, the basic operations used in the straight-line program preserve compatibility:
\begin{itemize}
   \item adding low-degree polynomials (which only affects low coefficients);
   \item multiplying by a \emph{known-power gadget} $H_{2^i}(x)$ (a fixed monic degree-$2^i$ polynomial, treated as known to the decoder),
         which shifts blocks of coefficients by $2^i$ positions without destroying recoverability of a top window;
   \item combining two compatible pairs in a block-triangular way so that the coefficient windows do not interfere.
\end{itemize}
These closure properties are what lets us scale the warm-up decoding argument beyond tiny degrees.

\paragraph{Recursion: filling coefficients with a BRW-style structure.}
At a high level, we build a decodable family $Q_{2^k-1}(x;\alpha)$ of monic degree $(2^k-1)$ polynomials using a divide-and-conquer pattern reminiscent of Rabin--Winograd / Bernstein constructions~\cite{rabin1972number}:
we keep a ladder of known-power gadgets $H_2,H_4,\dots,H_{2^{k-1}}$ (computed once and reused),
and we recursively ``fill'' disjoint coefficient windows by multiplying previously constructed pieces by the appropriate $H_{2^i}$ and adding a fresh low-degree correction.
The compatibility invariant guarantees that each recursive layer introduces new parameters that remain decodable from the top coefficients, despite coefficient mixing in intermediate multiplications.
The resulting circuits use $\Theta(2^{k-1})$ multiplications for degree $\Theta(2^k)$, i.e.\ about half the multiplications of Horner.

\paragraph{From the recursion to \cref{thm:main}.}
The recursion above yields splittable pairs for essentially all odd degrees (with a small number of explicit base cases).
Given a splittable pair $(T^{(1)},T^{(2)})$ for degree $n$, the circuit for $P(x)=xT^{(1)}(x)+T^{(2)}(x)$ uses one more multiplication than the shared computation of $T^{(1)}$ and $T^{(2)}$.
For even degrees, we reduce to the odd case by writing
\(
   P(x)=c_0 + x\cdot Q(x)
\)
where $Q$ is a monic degree-$(n-1)$ polynomial with coefficients $(c_1,\dots,c_{n-1})$; this costs one extra multiplication by $x$.
In all cases we obtain $\lfloor n/2\rfloor+1$ multiplications overall, and the decoding invariants ensure that the induced coefficient map is generically invertible by rational functions, i.e.\ admits rational preprocessing in the sense of \cref{sec:model}.

\subsection{\texorpdfstring{$k$}{k}-wise independent hashing and why bijectivity matters}
\label{sec:main-theorem:kwise}

\paragraph{Background: $k$-wise independence via polynomials.}
Let $\F$ be a field and fix $k\ge 1$.
A standard way to build a $k$-wise independent hash family $\F\to\F$ is to choose a random
degree-$k$ \emph{monic} polynomial and evaluate it on the input~\cite{wegman1981new}.
Concretely, for a coefficient vector $c=(c_0,\dots,c_{k-1})\in\F^k$ define
\begin{equation}
  h_c(x) \coloneqq x^k + c_{k-1}x^{k-1} + \cdots + c_1 x + c_0.
  \label{eq:kwise:monic-poly}
\end{equation}
If $c\gets\F^k$ is uniform, then the family $\{h_c\}$ is $k$-wise independent:
for any $k$ distinct inputs $x_1,\dots,x_k\in\F$, the random vector
$\bigl(h_c(x_1),\dots,h_c(x_k)\bigr)$ is uniform in $\F^k$.
Indeed,
\[
  h_c(x_j) - x_j^k = \sum_{i=0}^{k-1} c_i x_j^i,
\]
and the linear map $c\mapsto \bigl(\sum_i c_i x_j^i\bigr)_{j=1}^k$ is a Vandermonde transform,
hence invertible when the $x_j$ are distinct.

\paragraph{Why multiplications dominate in practice.}
Evaluating \eqref{eq:kwise:monic-poly} is a tight inner loop in many randomized data structures
(hash tables, sketches, filters) and in cryptographic-style workloads.
In typical finite-field implementations, additions are cheap but multiplications are expensive
because they entail modular reduction:
either integer multiplication plus reduction in $\F_p$, or carryless multiplication plus
reduction in $\F_{2^w}$.
Thus the main performance objective is often to reduce the number of field multiplications,
especially on the \emph{critical path} (dependent multiplications).

\paragraph{Classical evaluation baselines.}
The standard baseline is Horner’s rule, which evaluates \eqref{eq:kwise:monic-poly} using $k-1$
multiplications and $k$ additions:
\[
  \bigl(\cdots((x+c_{k-1})x + c_{k-2})x + \cdots \bigr)x + c_0.
\]
There are two widely used alternatives:
\begin{itemize}
  \item \textbf{Estrin / Paterson--Stockmeyer style.}
  These reorganize evaluation to expose instruction-level parallelism by grouping terms,
  reducing depth to $O(\log k)$ at the cost of extra multiplications (to form powers like $x^2,x^4,\dots$).
  This can help on wide machines when throughput matters more than total work~\cite{estrin1960organization,paterson1973evaluation}.
  \item \textbf{Rabin--Winograd preprocessing.}
  Rabin and Winograd~\cite{rabin1972number} showed that, with preprocessing of the coefficient
  vector $c$ using only rational operations, one can evaluate a degree-$k$ polynomial with
  about $k/2 + O(\log k)$ multiplications.
  Their work framed preprocessing in a way that is compatible with finite fields, and they
  conjectured their overhead beyond $k/2$ was essentially necessary.
\end{itemize}

\paragraph{Our setting: preprocessing as a parameterization.}
In this paper, we consider an evaluation circuit $E(x;\theta)$ whose \emph{runtime key} is a
parameter vector $\theta\in\F^d$.
For each $\theta$, the function $x\mapsto E(x;\theta)$ is a monic degree-$k$ polynomial, with
coefficient vector $\mathsf{coeff}(\theta)\in\F^k$ as in \cref{sec:model}.
There are then two natural ways to generate a $k$-wise hash function:
\begin{enumerate}
  \item sample coefficients $c\gets\F^k$ and preprocess them into parameters
        $\theta=\mathsf{pre}(c)$, then hash by $x\mapsto E(x;\theta)$;
  \item sample parameters $\theta$ directly (for example uniformly from $\F^d$) and hash by
        $x\mapsto E(x;\theta)$.
\end{enumerate}

\paragraph{Why bijectivity/birationality is important for hashing.}
The distinction above is exactly where the “bijective/injective” viewpoint matters.
If $\mathsf{coeff}$ is not injective, then distinct parameters can represent the same polynomial.
This is not merely an aesthetic issue: sampling $\theta$ induces a potentially \emph{biased}
distribution on polynomials (some polynomials may have many preimages, others none),
which can break clean statements like “the hash is $k$-wise independent” if the key-generation
procedure is “pick a random $\theta$”.

Rational preprocessing, as defined in \cref{sec:model}, is a formal way to rule out this bias:
it requires a rational right-inverse $\mathsf{pre}$ of $\mathsf{coeff}$ on a dense set, i.e.\ that
$\mathsf{coeff}$ is generically one-to-one (birational onto its image).
Operationally this gives two benefits:
\begin{itemize}
  \item \textbf{Soundness of preprocessing.}
  Given a uniformly random coefficient vector $c$ (the standard $k$-wise key), we can compute
  parameters $\theta=\mathsf{pre}(c)$ such that $E(x;\theta)\equiv h_c(x)$.
  This lets us implement $k$-wise hashing using the fast circuit while preserving the exact
  algebraic distribution that the usual Vandermonde proof assumes.
  \item \textbf{Optional ``sample-$\theta$'' viewpoint.}
  When $\mathsf{coeff}$ is generically injective, there is essentially one parameter vector per
  polynomial (up to a lower-dimensional exceptional set), so treating $\theta$ itself as the key
  corresponds to selecting a polynomial from a well-defined family without duplicates.
\end{itemize}
This is the sense in which “fast evaluation with rational preprocessing” becomes a
bijective/injective parameterization problem.

\paragraph{Implication of \cref{thm:main} for $k$-wise hashing.}
Applying \cref{thm:main} with $n=k$, we obtain an evaluation scheme for \eqref{eq:kwise:monic-poly}
that uses only $\lfloor k/2\rfloor+1$ multiplications after preprocessing the coefficients into
runtime parameters.
This improves over Horner’s $k-1$ multiplications and removes the $O(\log k)$ overhead in
Rabin--Winograd while staying within the rational preprocessing model.
We evaluate practical performance trade-offs against Horner, Estrin-style evaluation, and
Rabin--Winograd in \cref{sec:experiments:kwise}.
